\documentclass[11pt]{article}

\usepackage[margin=1in]{geometry}
\usepackage{amsmath,amssymb}
\usepackage{graphicx}
\usepackage{booktabs}
\usepackage{xcolor}
\usepackage[colorlinks=true,linkcolor=blue,citecolor=blue,urlcolor=blue]{hyperref}

\newcommand{\code}[1]{\texttt{#1}}
\newcommand{\pf}{p_{\mathrm{f}}}
\newcommand{\pfz}{p_{\mathrm{f}0}}
\newcommand{\pwh}{p_{\mathrm{wh}}}

\title{Referencing the Habanero~4 wellhead record to the initial fault
       pressure}
\author{Cooper Basin / HBI cycle-2 modelling}
\date{}

\begin{document}
\maketitle

\begin{abstract}
\noindent
The Habanero~4 pressure record is an \emph{absolute wellhead} pressure, while
HBI computes a \emph{change in fault pore pressure} from $\pf = 0$. Comparing
them requires the record to be carried down the wellbore and referenced to the
fault's initial pore pressure $\pfz$. This note derives that conversion,
assembles every term from reported field measurements, and states the resulting
shift. Referencing to $\pfz$ rather than to the first sample of the wellhead
record raises the observed pressure change by $2.78$~MPa at the start of the
modelled window, easing to $1.53$~MPa by day~10 as the rate-dependent pipe
friction grows. No simulation input is altered; this is a correction to how the
\emph{data} are reduced.
\end{abstract}

\section{The quantity each side computes}

The field measurement is $\pwh(t)$, an absolute pressure at the wellhead gauge.
HBI initialises its pore-pressure perturbation to zero (\code{pfinit = 0}) and
its effective normal stress \code{sigmainit} already contains the initial fault
pore pressure, so the model's output is
\begin{equation}
  \Delta p_{\mathrm{sim}}(t) \;=\; \pf(t) - \pfz .
\end{equation}
To place the measurement on the same footing, the wellhead pressure must be
carried to the fault and referenced to the same $\pfz$:
\begin{equation}
  \pf(t) \;=\; \pwh(t) \;+\; \rho g Z \;-\; \Delta p_{\mathrm{fric}}\!\left(q(t)\right),
  \label{eq:carry}
\end{equation}
\begin{equation}
  \boxed{\;\Delta p_{\mathrm{obs}}(t) \;=\; \pwh(t)
         \;-\; \Delta p_{\mathrm{fric}}\!\left(q(t)\right)
         \;-\; \underbrace{\left[\,\pfz - \rho g Z\,\right]}_{\textstyle \text{datum}}\;}
  \label{eq:dp}
\end{equation}
where $Z$ is the vertical distance from the gauge to the fault, $\rho$ the mean
density of the fluid column, and $\Delta p_{\mathrm{fric}}$ the frictional loss
along the flow path.

\paragraph{The hydrostatic head does not cancel.}
It is tempting to assume the $\rho g Z$ term drops out of a difference. It does
so only if the reference is \emph{itself} defined as a wellhead pressure. Here
the reference is derived from $\pfz$, a pressure at depth, so as
Eq.~\eqref{eq:dp} shows $\rho g Z$ enters once and additively. Over
$Z \approx 4.2$~km this makes the answer directly sensitive to $\rho$, at
\begin{equation}
  \frac{\partial (\text{datum})}{\partial \rho} \, \delta\rho
  \;=\; -g Z \, \delta\rho
  \;=\; -2.05~\mathrm{MPa} \quad \text{per } 50~\mathrm{kg\,m^{-3}} .
  \label{eq:sens}
\end{equation}
Every term in Eq.~\eqref{eq:dp} is therefore assembled below from a reported
measurement rather than assumed.

\section{The initial fault pressure $\pfz$}

Holl \& Barton \cite{holl2015} report a reservoir pressure of
\begin{equation}
  \pfz = 72.70~\mathrm{MPa} \quad \text{at } -4100~\mathrm{m\,AHD},
\end{equation}
the median depth of the Habanero Fault, corresponding to roughly $34$~MPa of
overpressure above hydrostatic. Hogarth \& Bour \cite{hogarth2015} quote
$73.0$~MPa at $-4140$~m~BSL for the same reservoir, consistent to $0.3$~MPa
once the $40$~m datum difference is allowed for.

\paragraph{A caution on circularity.}
The model's effective normal stress, \code{sigmainit} $=27.99$~MPa, is
\emph{not} an independent check on $\pfz$. Resolving the measured stress state
($S_v = 100$, $S_{H\max} = 160$~MPa, fault dip $10^\circ$) gives
$\sigma_{nn} = 101.81$~MPa, so $\sigma_{nn} - \code{sigmainit} = 73.82$~MPa ---
but that is the value used by Wang \& Dunham \cite{wang2022}, from which
\code{sigmainit} was itself derived. The agreement is therefore definitional,
not evidential. Note also that adopting Holl's $72.70$~MPa implies an effective
normal stress of $29.11$~MPa rather than $27.99$; the two differ by
$1.12$~MPa and cannot be varied independently, since
$\tau_0 = \mu_0\,\code{sigmainit}$.

\section{The column length $Z$}

Holl's datum is referenced to AHD (approximately sea level), whereas the
pressure gauge sits at the wellhead. The Habanero~4 Well Completion Report
\cite{wcr} gives a ground level of $73.34$~m~AHD and a rotary-table-to-ground
distance of $9.14$~m, placing the gauge $82.48$~m \emph{above} the datum.
Hence
\begin{equation}
  Z = 4100.0 + 73.34 + 9.14 = 4182.5~\mathrm{m},
\end{equation}
$82.5$~m longer than the fault depth alone. By Eq.~\eqref{eq:sens} that is
worth $0.81$~MPa, all of it lowering the datum and so raising
$\Delta p_{\mathrm{obs}}$.

\section{The column density $\rho$: two fluid states}

The wellbore does not contain one fluid. Before injection the column had been
standing in a reservoir at $240$--$250^\circ$C and had warmed towards the
geothermal profile: hot and light. During injection it carries surface water at
$25$--$50$~L\,s$^{-1}$, with a transit time of order twenty minutes and
correspondingly little opportunity to heat: cold and dense. A single $\rho$
cannot represent both, and by Eq.~\eqref{eq:sens} the difference is not small.

\paragraph{Flowing column (the one that matters).}
The WCR \cite{wcr} records the completion fluid left in the well as
$8.35$~ppg ambient, i.e.
\begin{equation}
  \rho_{\mathrm{flow}} = 1001~\mathrm{kg\,m^{-3}},
\end{equation}
which is the wellbore column immediately before the October/November~2012
stimulation. Independently, Hogarth \& Bour \cite{hogarth2015} Table~3 measures
the cold limb of the 2013 closed loop at $983$--$1006$~kg\,m$^{-3}$ for
$80^\circ$C brine. The value used is therefore reported, not selected.

\paragraph{Static column (a consistency check).}
Inverting Eq.~\eqref{eq:dp} at $t<0$, when the well was shut and in equilibrium
with the reservoir, the measured pre-injection wellhead pressure of
$33.970$~MPa requires
\begin{equation}
  \rho_{\mathrm{static}}
  = \frac{\left(\pfz - \pwh^{\,\mathrm{static}}\right)}{g Z}
  = \frac{(72.70 - 33.970)\times 10^{6}}{9.81 \times 4182.5}
  = 944~\mathrm{kg\,m^{-3}},
\end{equation}
consistent with a mean column temperature near $90$--$100^\circ$C, which is
what a shut-in well in this reservoir should exhibit. Hogarth \& Bour's
independently measured static surface pressure of $33.7$~MPa gives $949$. Had
$\pfz$ been badly wrong, this inversion would have returned an unphysical
density; that it does not is the check.

The head difference between the two states,
$(1001-944)\,g Z = 2.34$~MPa, is the physical origin of the shift reported
here. It is not an assumption layered on top of the correction: it \emph{is}
the correction.

\section{Frictional loss $\Delta p_{\mathrm{fric}}$}

The flow path is not a single pipe, and the diameter previously used was the
wrong one. The Darcy--Weisbach loss over a segment of length $L$ and internal
diameter $D$ at volumetric rate $q$ is
\begin{equation}
  \Delta p_{\mathrm{fric}} = f\,\frac{8 L \rho q^{2}}{\pi^{2} D^{5}},
  \qquad f = 0.015 .
  \label{eq:fric}
\end{equation}
The $D^{-5}$ dependence makes the choice of diameter decisive. The WCR
\cite{wcr} completion tally gives $7''$ $41\#$ tubing of \emph{internal}
diameter $5.820'' = 0.14783$~m run to the production packer at
$3005.37$~m~MDRT, with wider $9\text{-}7/8''$ casing below. Earlier work used
$D = 0.1778$~m uniformly, which is the tubing \emph{outside} diameter, over the
whole $4077$~m. Summing Eq.~\eqref{eq:fric} over the two real segments:

\begin{center}
\begin{tabular}{lrrr}
  \toprule
  & $L$ (m) & $D$ (m) & $\Delta p_{\mathrm{fric}}$ at $60.9$~L\,s$^{-1}$ (MPa) \\
  \midrule
  tubing, $5.820''$ ID   & 3005.0 & 0.14783 & 1.93 \\
  casing, $8.625''$ ID   & 1177.5 & 0.21908 & 0.10 \\
  \midrule
  total                  &        &         & \textbf{2.03} \\
  previous, uniform $0.1778$~m & 4077.0 & 0.17780 & 1.03 \\
  \bottomrule
\end{tabular}
\end{center}

\noindent
a factor of $1.97$, essentially all of it in the tubing. Because
$\Delta p_{\mathrm{fric}}$ enters Eq.~\eqref{eq:dp} directly, the previous
value was a $1$~MPa error in the observed pressure change at peak rate. It
scales as $q^{2}$, so it matters least on the $4$--$8$~d plateau
($q \approx 23$~L\,s$^{-1}$, $0.28$~MPa) and most at the day-2.5 rate peak.

\section{The datum, and the resulting shift}

Collecting the terms,
\begin{equation}
  \text{datum} = \pfz - \rho_{\mathrm{flow}} g Z
               = 72.70 - 41.071
               = \boxed{31.629~\mathrm{MPa}}
\end{equation}
expressed at the wellhead gauge. Table~\ref{tab:series} gives the resulting
series against the previous convention, in which the record's first sample
($34.412$~MPa, at data-day~0) was subtracted. Simulation time is measured from
data-day~$4.300$, where injection resumes.

\begin{table}[htbp]
\centering
\caption{Measured pressure change on the two conventions. ``previous''
subtracts the record's first sample; ``this note'' applies
Eq.~\eqref{eq:dp}. The shift is not constant because
$\Delta p_{\mathrm{fric}} \propto q^{2}$.}
\label{tab:series}
\begin{tabular}{rrrrr}
  \toprule
  $t_{\mathrm{sim}}$ (d) & $\pwh$ (MPa) & previous (MPa) & this note (MPa) & shift (MPa) \\
  \midrule
   0.00 & 33.274 & $-1.138$ & $+1.645$ & $+2.783$ \\
   1.00 & 34.184 & $-0.228$ & $+2.367$ & $+2.595$ \\
   2.00 & 39.307 & $+4.895$ & $+7.214$ & $+2.319$ \\
   4.00 & 43.492 & $+9.080$ & $+11.472$ & $+2.392$ \\
   6.00 & 43.437 & $+9.025$ & $+11.524$ & $+2.499$ \\
   8.00 & 43.430 & $+9.018$ & $+11.517$ & $+2.499$ \\
   8.70 & 45.112 & $+10.700$ & $+12.732$ & $+2.031$ \\
  10.00 & 52.035 & $+17.623$ & $+19.149$ & $+1.526$ \\
  12.00 & 48.608 & $+14.196$ & $+16.227$ & $+2.031$ \\
  \bottomrule
\end{tabular}
\end{table}

\paragraph{The fault begins above virgin pressure.}
At $t_{\mathrm{sim}} = 0$ the measured $\Delta p$ is $+1.65$~MPa, not zero: the
fault is above its virgin pressure when injection resumes, notwithstanding that
the well had been vented during the preceding shut-in. This is the residual
formation overpressure left by the first injection cycle, and it is why
referencing to $\pfz$ \emph{raises} the curve rather than lowering it. A
convention that forces $\Delta p(0) = 0$ discards this information.

\section{What remains assumed}

\begin{itemize}
\item \textbf{$\rho$ is a depth-averaged constant.} The real column has a
      temperature profile, so $\rho$ varies with depth and, during injection,
      with time as the cold front advances down the tubing. The WCR contains a
      full-well temperature log (as a vector plot requiring digitisation) and a
      Core Labs PVT study with lab-measured brine density
      ($0.9193$~g\,cm$^{-3}$ at $200^\circ$C, $0.9990$ at $100^\circ$C, both at
      $10\,600$~psig); either would replace the constant with a profile.
\item \textbf{$f = 0.015$} is inherited, not measured. The WCR reports it as
      determined ``from trial and error''.
\item \textbf{$\pfz$ and \code{sigmainit} are mutually inconsistent by
      $1.12$~MPa}, as noted in \S2. Resolving this requires moving both
      $\pfz$ and $\tau_0$ together.
\item \textbf{Two-phase or gas effects} in the column are neglected entirely.
\end{itemize}

\section*{Reproducibility}

Every number above is computed by
\code{docs/cooper\_basin/fault\_pressure.py}, which parses the reported values
and writes \code{docs/figs/cycle2/DATUM.md} together with the figure
\code{docs/figs/cycle2/fault\_pressure.pdf}. The scoring script
(\code{score\_cycle2.py}) and the comparison figures
(\code{compare\_arms.py}, \code{compare\_disc.py}) import the datum from that
module rather than restating it.

\begin{thebibliography}{9}

\bibitem{holl2015}
H.-G.~Holl and C.~Barton,
\emph{Habanero Field --- Structure and State of Stress},
Proc.\ World Geothermal Congress 2015, paper 31034.
\url{https://pangea.stanford.edu/ERE/db/WGC/papers/WGC/2015/31034.pdf}

\bibitem{hogarth2015}
R.~A.~Hogarth and D.~Bour,
\emph{Flow Performance of the Habanero EGS Closed Loop},
Proc.\ World Geothermal Congress 2015, paper 31006.
\url{https://pangea.stanford.edu/ERE/db/WGC/papers/WGC/2015/31006.pdf}

\bibitem{wcr}
F.~Walsh (Geodynamics Ltd),
\emph{Well Completion Report, Habanero-4},
HPP-FN-OT-RPT-00492-1.0, 8 January 2013. PEPS South Australia well~2739,
open file.
\url{https://onepeps-api.azurewebsites.net/api/file/WellCompletionReport/2739/Habanero4_report.pdf}

\bibitem{wang2022}
T.~Wang and E.~M.~Dunham,
\emph{Hydraulic stimulation and fault slip in the Cooper Basin},
Scientific Reports \textbf{12}, 2022.

\end{thebibliography}

\end{document}
